\section{Two lemmas and a computable step}
The two lemmas in this section use \eqref{eq:gap_smoothness}
but do not require convexity or quasar-convexity of $f$.

\begin{lemma}[Transport and an upper model]\label{lem:transport}
Let $s,d$ be given by \eqref{eq:gap}. For $0\le t\le1$,
\begin{align}
 G(x+td)&\le(1-t)G(x)
    +D\|d\|\int_0^t\bigl(L_0+L_1G(x+ud)\bigr)\,du,
       \label{eq:transport}\\
 L_0+L_1G(x+td)&\le
      \bigl(L_0+L_1G(x)\bigr)e^{L_1D\|d\|t}.
       \label{eq:transport_exp}
\end{align}
Consequently, for $0\le\alpha\le1$,
\begin{equation}\label{eq:upper_model}
 f(x+\alpha d)\le f(x)-\alpha G(x)
 +(L_0+L_1G(x))\|d\|^2
       \int_0^\alpha(\alpha-t)e^{L_1D\|d\|t}\,dt.
\end{equation}
\end{lemma}
\begin{proof}
For every $v\in Q$, expand the inner product at $x$:
\[
 \langle\nabla f(x+td),x+td-v\rangle
 \le G(x)-tG(x)
       +D\|\nabla f(x+td)-\nabla f(x)\|.
\]
Maximizing over $v$ and integrating the Hessian along the segment
gives \eqref{eq:transport}. Drop the nonpositive term $-tG(x)$
and apply Gronwall's inequality to $L_0+L_1G(x+td)$ to obtain
\eqref{eq:transport_exp}; the case $L_1=0$ is immediate.
Taylor's formula with integral remainder gives \eqref{eq:upper_model},
since
$|d^\top\nabla^2f(x+td)d|\le\|\nabla^2f(x+td)\|\|d\|^2$.
\end{proof}

The exponential in \eqref{eq:upper_model} bounds the growth of
curvature while moving. The square $\|d\|^2$ comes from the second
directional derivative. Both arise in the proof; neither can be
removed by substituting the gap for a gradient norm in another model.
The selected segment is essential. For example,
$f(t)=t+t^3/6$ on $[0,1]$ has $|f''(t)|\le G(t)$, yet $G(0)=0$
and $f(1)>f(0)+f'(0)$. Thus an upper model with initial coefficient
$L_0+L_1G(x)$ cannot hold for \emph{every} pair $x,y$.

\paragraph{Step rule.}
If $G(x)>0$ and $L_1>0$, minimize the right-hand side of
\eqref{eq:upper_model} over $[0,1]$:
\begin{equation}\label{eq:log_step}
 \boxed{\displaystyle
 \alpha=\min\left\{1,\,
 \frac{1}{L_1D\|d\|}
 \log\left(1+
 \frac{L_1D\,G(x)}
 {(L_0+L_1G(x))\|d\|}\right)\right\}.}
\end{equation}
Here and below $\log$ is the natural logarithm.
If $L_1=0<L_0$, use the continuous limit
$\alpha=\min\{1,G(x)/(L_0\|d\|^2)\}$.
If $L_0=L_1=0$, use $\alpha=1$; $f$ is then affine on $Q$.
If $G(x)=0$, stop before evaluating any quotient.

\begin{lemma}[Descent]\label{lem:descent}
The step above satisfies
\begin{equation}\label{eq:descent}
 f(x+\alpha d)\le f(x)-\frac{\alpha G(x)}2.
\end{equation}
Outside the affine case,
\begin{equation}\label{eq:step_lower}
 \alpha\ge
 \min\left\{1,
   \frac{(\log2)G(x)}{(L_0+L_1G(x))D^2}\right\}.
\end{equation}
\end{lemma}
\begin{proof}
For this scalar calculation put $b=L_1D\|d\|>0$.
The integral in \eqref{eq:upper_model} equals
$(e^{b\alpha}-1-b\alpha)/b^2$.
Differentiating the model gives
\[
 -G(x)+(L_0+L_1G(x))\|d\|^2
             \frac{e^{b\alpha}-1}{b}.
\]
Its zero gives \eqref{eq:log_step}. At the capped minimizer
the derivative is nonpositive. Since $(e^{bt}-1)/b$ is convex
and vanishes at zero, its integral on $[0,\alpha]$ is at most
half of $\alpha$ times its endpoint value. Substitution in
\eqref{eq:upper_model} proves \eqref{eq:descent}.

For \eqref{eq:step_lower}, temporarily set
$r=\|d\|/D\le1$ and
$q=L_1G(x)/(L_0+L_1G(x))\le1$. The uncapped step is
$\log(1+q/r)/(L_1D^2r)$, and
\[
 \log(1+q/r)\ge\log(1+q)\ge q\log2\ge qr\log2.
\]
The middle inequality is concavity of the logarithm.
When $L_1=0$, both assertions follow from the quadratic model.
\end{proof}

In the norm-based step of~\cite{vyguzov}, a restriction on the
step length bounds an exponential by $e$.
Here \eqref{eq:log_step} solves the exponential model directly.
In particular, its logarithm makes
\[
 e^{L_1D\|d\|\alpha}-1
 \le\frac{L_1D\,G(x)}{(L_0+L_1G(x))\|d\|},
\]
which is precisely the inequality used in the descent proof.

\paragraph{Algorithm and stopping.}
Fix an accuracy $\varepsilon>0$.
At each visited point $x_k$, evaluate $f(x_k)$, its gradient,
and a full linear minimizer. Retain
\begin{equation}\label{eq:stop}
 \ell_k=\max_{0\le j\le k}
        \left\{f(x_j)-\frac{G(x_j)}{\gamma}\right\},
 \qquad \text{stop if } f(x_k)-\ell_k\le\varepsilon.
\end{equation}
Otherwise set $x_{k+1}=x_k+\alpha_kd_k$ using \eqref{eq:log_step}
or its stated limiting cases. Each new function value is evaluated
once. There is no line search.
